\documentclass{standalone}
\usepackage{circuitikz}
\usetikzlibrary{calc}
\definecolor{circuitnotemark}{HTML}{FE00FE}
\makeatletter
\pgfdeclareshape{smacoax}{
  \anchor{center}{\pgfpointorigin}
  \anchor{text}{\pgfpointorigin}
  \anchor{north}{\pgfpoint{0cm}{0.3cm}}
  \anchor{south}{\pgfpoint{0cm}{-0.3cm}}
  \anchor{east}{\pgfpoint{0.3cm}{0cm}}
  \anchor{west}{\pgfpoint{-0.3cm}{0cm}}
  \anchor{pin 1}{\pgfpoint{-0.7cm}{0cm}}
  \anchor{pin 2}{\pgfpoint{0cm}{-0.7cm}}
  \anchor{bpin 1}{\pgfpoint{-0.3cm}{0cm}}
  \anchor{bpin 2}{\pgfpoint{0cm}{-0.3cm}}
  \backgroundpath{
    \pgfpathmoveto{\pgfpoint{-0.274cm}{-0.122cm}}
    \pgfpatharc{204}{516}{0.3cm}
    \pgfpathmoveto{\pgfpoint{-0.7cm}{0cm}}\pgfpathlineto{\pgfpointorigin}
    \pgfpathmoveto{\pgfpoint{0cm}{-0.7cm}}\pgfpathlineto{\pgfpoint{0cm}{-0.3cm}}
  }
  \foregroundpath{
    \pgfpathcircle{\pgfpointorigin}{0.07cm}
    \pgfusepath{fill}
  }
}
\makeatother
\begin{document}
\begin{circuitikz}[american, line width=0.8pt]
\ctikzset{bipoles/length=1.2cm}
\ctikzset{grounds/scale=1.36}
\coordinate (b2) at (2,-2);
\coordinate (c5) at (8,-4);
\coordinate (e5) at (8,-8);
\coordinate (b8) at (14,-2);
\coordinate (c2) at (2,-4);
\coordinate (c8) at (14,-4);
\coordinate (b5) at (8,-2);
\node[smacoax, draw, xscale=-1] (part-J1) at (b2) {}; % line 3
\node[anchor=south] at (part-J1.north) {$J_{1}$}; % line 3
\draw (c5) to[TL, l_=$S_{1}$, a^=$50\,\Omega$] (e5); % line 4
\node[smacoax, draw] (part-J2) at (b8) {}; % line 5
\node[anchor=south] at (part-J2.north) {$J_{2}$}; % line 5
\node[ground] at (c2) {}; % line 6
\node[ground] at (c8) {}; % line 7
\draw (part-J1.pin 1) -- (b5); % line 9
\draw (b5) -- (part-J2.pin 1); % line 9
\draw (b5) -- (c5); % line 10
\draw (part-J1.pin 2) -- (c2); % line 11
\draw (part-J2.pin 2) -- (c8); % line 12
\node[circ] at (b5) {};
\path (1.619,-0.593) rectangle (2.381,0.593); % line 14
\node[anchor=west, inner sep=0, circuitnotemark, font=\large] at (2,0) {X}; % line 14
\path (13.619,-0.593) rectangle (14.381,0.593); % line 15
\node[anchor=west, inner sep=0, circuitnotemark, font=\large] at (14,0) {X}; % line 15
\path (7.556,-9.193) rectangle (8.445,-8.007); % line 16
\node[anchor=west, inner sep=0, circuitnotemark, font=\large] at (8,-8.6) {X}; % line 16
\node[anchor=south west, inner sep=2pt, circuitnotemark, font=\LARGE\bfseries] (circuittitle) at (current bounding box.north west) {X};
\path (circuittitle.south west) rectangle ++(12.131,0.853);
\end{circuitikz}
\end{document}
